\documentclass[11pt]{article}
\usepackage[margin=0.7in]{geometry}
\usepackage{amsmath,amssymb}
\usepackage{enumitem}
\usepackage{tikz}
\usepackage{pgfplots}
\pgfplotsset{compat=1.18}
\usepackage{fancyhdr}
\usepackage{parskip}
\setlist[enumerate]{leftmargin=*, itemsep=6pt}
\pagestyle{fancy}
\fancyhf{}
\fancyfoot[L]{\small\textit{en}Vision\textsuperscript{\textregistered} Algebra 2}
\fancyfoot[C]{\small 4-5 Lesson Quiz, short form}
\fancyfoot[R]{\small\thepage}
\renewcommand{\headrulewidth}{0pt}
\newcommand{\choice}[1]{\tikz[baseline=-0.6ex]\node[draw, circle, inner sep=0.6pt, minimum size=3.3mm, font=\scriptsize\bfseries] {#1};\;}
\newcommand{\checkbox}{$\square$\;}
\pgfplotsset{quizgrid/.style={
  axis lines=middle, grid=both, grid style={gray!45},
  ticklabel style={font=\scriptsize, fill=white, inner sep=1pt},
  xlabel={$x$}, ylabel={$y$}, axis line style={thick}, clip=true}}
\newcommand{\worklines}[1]{\foreach \i in {1,...,#1}{\vspace{16pt}\rule{\linewidth}{0.4pt}\par}}

\begin{document}

\noindent\textbf{Name}\; \rule{6cm}{0.4pt} \hfill Period \rule{1.2cm}{0.4pt} \quad Date \rule{2.2cm}{0.4pt}

\vspace{4pt}\noindent\rule{\textwidth}{0.6pt}

\vspace{2pt}\noindent{\LARGE\textbf{4-5 \; Lesson Quiz: Solving Rational Equations}} \hfill {\small short form}

\vspace{4pt}\noindent\rule{\textwidth}{0.6pt}

{\small Question numbers match the Savvas Lesson Quiz. Questions 1 and 5 are not on this form. Show your work for every question: what you multiplied by, each step, and your check.}

\vspace{4pt}
\begin{minipage}[t]{0.46\textwidth}
\textbf{2.} Aaron can paint a door in 6 hours, and Anna can do the same job in 3 hours. Working together, how many hours will it take them to paint a door?

\vspace{8pt}
Equation: \rule{4cm}{0.4pt}

\vspace{10pt}
Answer: $h=$ \rule{2cm}{0.4pt} hours
\end{minipage}\hfill
\begin{minipage}[t]{0.50\textwidth}
\textbf{Work:}\par
\worklines{4}
\end{minipage}

\vspace{12pt}\noindent\rule{\textwidth}{0.4pt}

\vspace{6pt}
\begin{minipage}[t]{0.46\textwidth}
\textbf{3.} Solve the equation $\dfrac{3}{x-1}=\dfrac{4}{2x+3}$.

\vspace{8pt}
Check in the original equation: \rule{2.6cm}{0.4pt}

\vspace{10pt}
Solution: $x=$ \rule{2.4cm}{0.4pt}
\end{minipage}\hfill
\begin{minipage}[t]{0.50\textwidth}
\textbf{Work:}\par
\worklines{4}
\end{minipage}

\vspace{12pt}\noindent\rule{\textwidth}{0.4pt}

\vspace{6pt}
\begin{minipage}[t]{0.46\textwidth}
\textbf{4.} Solve the equation $\dfrac{-3}{x+3}=\dfrac{x}{x+3}-\dfrac{x}{5}$. Identify any extraneous solutions.

\vspace{8pt}
Solution: \rule{3.4cm}{0.4pt}

\vspace{10pt}
Extraneous solution: \rule{2.6cm}{0.4pt}

\vspace{10pt}
Why is it extraneous?

\vspace{14pt}\rule{\linewidth}{0.4pt}
\end{minipage}\hfill
\begin{minipage}[t]{0.50\textwidth}
\textbf{Work:}\par
\worklines{5}
\end{minipage}

\vfill
\noindent\textbf{Scoring:} Question 3 and Question 4 are 10 points each. Question 2 is a bonus of up to 5 points. Total out of 20.

\newpage
\noindent{\LARGE\textbf{4-5 Lesson Quiz, short form: answer key}}

\vspace{4pt}\noindent\rule{\textwidth}{0.6pt}

\begin{enumerate}
\item[\textbf{2.}] \textbf{$h=2$ hours.} Aaron paints $\tfrac16$ of the door per hour and Anna $\tfrac13$, so $\tfrac16+\tfrac13=\tfrac1h$. Multiply by $6h$: $h+2h=6$, so $3h=6$ and $h=2$. Sense check: 2 hours is less than either time alone. Savvas quiz item 2, DOK 2, skill A87; taught by Example 2 (cut in the fast packet), rehearsed by the packet's Practice 2 (TE Practice 26, Kenji and Oscar). \textbf{Bonus 5:} 3 for the rate equation $\tfrac16+\tfrac13=\tfrac1h$ (or an equivalent), 2 for $h=2$. An answer of 4.5 hours (the average of 6 and 3) or 9 hours (the sum) earns 0.
\item[\textbf{3.}] \textbf{$x=-6.5$} $\left(-\tfrac{13}{2}\right)$. Domain: $x\neq1$ and $x\neq-\tfrac32$. Multiply both sides by $(x-1)(2x+3)$: $3(2x+3)=4(x-1)$, so $6x+9=4x-4$, $2x=-13$, $x=-6.5$. Check: $\tfrac{3}{-7.5}=-0.4$ and $\tfrac{4}{-10}=-0.4$ \checkmark; $-6.5$ is in the domain. Savvas quiz item 3, DOK 2, skill A87; taught by Example 1 (TE Practice 32, the same structure, is tagged SEE EXAMPLE 1); no kept Example or Try It shows a variable denominator on each side, so the packet's Practice 1 is TE Practice 32 ($\tfrac{3}{x+1}=\tfrac{2}{x-3}$, answer D, $x=11$). \textbf{10 points:} 3 for multiplying by both denominators (or cross-multiplying), 3 for $6x+9=4x-4$, 2 for $x=-6.5$, 2 for the check (both sides equal $-0.4$).
\item[\textbf{4.}] \textbf{Solution: $x=5$; extraneous solution: $x=-3$.} Domain: $x\neq-3$. Multiply both sides by $5(x+3)$: $-15=5x-x(x+3)=5x-x^2-3x$, so $x^2-2x-15=0$, $(x-5)(x+3)=0$, $x=5$ or $x=-3$. $-3$ makes the denominator $x+3$ equal to 0, so it is not in the domain of the original equation. Check $x=5$: $\tfrac{-3}{8}$ and $\tfrac58-1=-\tfrac38$ \checkmark. Savvas quiz item 4, DOK 2, skill A87; taught by Example 4 (Part B, $\tfrac{3}{x-3}=\tfrac{x}{x-3}-\tfrac{x}{4}$, has the same structure). \textbf{10 points:} 3 for clearing the fractions with $5(x+3)$, 3 for $x^2-2x-15=0$ factored, 2 for $x=5$, 2 for naming $-3$ as extraneous with the reason (denominator 0 / not in the domain). Listing both $5$ and $-3$ as solutions earns at most 6.
\end{enumerate}

\textbf{Items left off this form} (available for a retake or Form B): Savvas item 1, solve $\frac{1}{2x-3}=4$ (choices A~$\frac12$, B~$\frac{8}{13}$, C~$\frac{11}{8}$, D~$\frac{13}{8}$; answer D, $\tfrac{13}{8}$; DOK 2, A87; Example 1, the same move as Try It 1); Savvas item 5, the boat in a stream: speed 10 mi/h in still water, 5 miles upstream and 5 miles back in 1 hour 20 minutes (answer D, 5 mi/h, from $\tfrac{5}{10-c}+\tfrac{5}{10+c}=\tfrac43$, $c^2=25$, the negative root rejected; DOK 2, A87; Example 5, cut in the fast packet).

\textbf{IXL replacement:} SmartScore 70 or above on CHP and CJC replaces this quiz with 20/20.

\end{document}
